% GENERATED FILE. Edit canonical lessons, then run npm run generate:books.
% canonical-source-hash: fnv1a64:9e9c712528ffa661
% canonical-lessons: UR-C72-ucha, UR-C72-nicha, UR-C72-bhara, UR-C72-khali, UR-C72-mitha

\chapter{High, Low, Full, Empty, Sweet}
\label{ch:ur-high-low-full-empty-sweet}
\hlchaptermodality{72}

\begin{chapteropening}
\textbf{By the end of this chapter:} \emph{I can say high, low, full, empty and sweet.}

Complete the last lesson of chapter 72: I can say high, low, full, empty and sweet.
\end{chapteropening}

\section[\texorpdfstring{\ur{اونچا} (ū̃chā)}{اونچا (ū̃chā)}]{\ur{اونچا} (ū̃chā) — high}

\label{lesson:UR-C72-ucha}



\begin{quote}
Before the new one: say the Urdu for expensive, then the Urdu for cheap.
\end{quote}

\subsection*{You'll want to know: \ur{اونچا}}
\textbf{\ur{اونچا}} — \emph{ū̃chā} — \textquotedblleft{}high\textquotedblright{}.

\begin{cousinweb}
\textbf{\ur{اونچا}} (\emph{ū̃chā}) is a tadbhava, a word inherited from Sanskrit and reshaped by sound change. It comes from \emph{ucca}, \textquotedblleft{}tall, lofty\textquotedblright{} (CDIAL 1634). CDIAL calls the influence of \emph{udañc}, \textquotedblleft{}turned upwards\textquotedblright{}, on the nasal very doubtful.
\end{cousinweb}

\begin{grammarlens}[title={What you've built}]
One more word for this chapter.
\end{grammarlens}

\subsection*{Your turn}
\begin{itemize}
\raggedright
  \item \emph{Say it:} \emph{ū̃chā}
  \item \emph{Say it:} \emph{ū̃chā}, once more
  \item \emph{Say it:} say \emph{ū̃chā}
  \item \emph{Recall:} say the Urdu for expensive, then the Urdu for cheap, then say \emph{ū̃chā} again
\end{itemize}

\begin{tcolorbox}[breakable,colback=teal!4,colframe=teal!35!black,title={Before you move on}]
What does \ur{اونچا} mean? (\textquotedblleft{}High\textquotedblright{}.) Say it once more.
\end{tcolorbox}
\section[\texorpdfstring{\ur{نیچا} (nīchā)}{نیچا (nīchā)}]{\ur{نیچا} (nīchā) — low}

\label{lesson:UR-C72-nicha}



\begin{quote}
Before the new one: say the Urdu for cheap, then the Urdu for high.
\end{quote}

\subsection*{You'll want to know: \ur{نیچا}}
\textbf{\ur{نیچا}} — \emph{nīchā} — \textquotedblleft{}low\textquotedblright{}.

\begin{cousinweb}
\textbf{\ur{نیچا}} (\emph{nīchā}) is a tadbhava, a word inherited from Sanskrit and reshaped by sound change. It comes from \emph{nīca}, \textquotedblleft{}low\textquotedblright{} (CDIAL 7540). CDIAL explains its \emph{-c-} as shaped after its opposite, \emph{ucca}, \textquotedblleft{}high\textquotedblright{}.
\end{cousinweb}

\begin{grammarlens}[title={What you've built}]
One more word for this chapter.
\end{grammarlens}

\subsection*{Your turn}
\begin{itemize}
\raggedright
  \item \emph{Say it:} \emph{nīchā}
  \item \emph{Say it:} \emph{nīchā}, once more
  \item \emph{Say it:} say \emph{nīchā}
  \item \emph{Recall:} say the Urdu for cheap, then the Urdu for high, then say \emph{nīchā} again
\end{itemize}

\begin{tcolorbox}[breakable,colback=teal!4,colframe=teal!35!black,title={Before you move on}]
What does \ur{نیچا} mean? (\textquotedblleft{}Low\textquotedblright{}.) Say it once more.
\end{tcolorbox}
\section[\texorpdfstring{\ur{بھرا} (bharā)}{بھرا (bharā)}]{\ur{بھرا} (bharā) — full}

\label{lesson:UR-C72-bhara}



\begin{quote}
Before the new one: say the Urdu for high, then the Urdu for low.
\end{quote}

\subsection*{You'll want to know: \ur{بھرا}}
\textbf{\ur{بھرا}} — \emph{bharā} — \textquotedblleft{}full\textquotedblright{}.

\begin{cousinweb}
\textbf{\ur{بھرا}} (\emph{bharā}) is a tadbhava, a word inherited from Sanskrit and reshaped by sound change. It comes from \emph{bharati}, \textquotedblleft{}bears, fills\textquotedblright{} (CDIAL 9397). \emph{bharā}, \textquotedblleft{}full\textquotedblright{}, is the past form of \emph{bharnā}, \textquotedblleft{}to fill\textquotedblright{}.
\end{cousinweb}

\begin{grammarlens}[title={What you've built}]
One more word for this chapter.
\end{grammarlens}

\subsection*{Your turn}
\begin{itemize}
\raggedright
  \item \emph{Say it:} \emph{bharā}
  \item \emph{Say it:} \emph{bharā}, once more
  \item \emph{Say it:} say \emph{bharā}
  \item \emph{Recall:} say the Urdu for high, then the Urdu for low, then say \emph{bharā} again
\end{itemize}

\begin{tcolorbox}[breakable,colback=teal!4,colframe=teal!35!black,title={Before you move on}]
What does \ur{بھرا} mean? (\textquotedblleft{}Full\textquotedblright{}.) Say it once more.
\end{tcolorbox}
\section[\texorpdfstring{\ur{خالی} (khālī)}{خالی (khālī)}]{\ur{خالی} (khālī) — empty}

\label{lesson:UR-C72-khali}



\begin{quote}
Before the new one: say the Urdu for low, then the Urdu for full.
\end{quote}

\subsection*{You'll want to know: \ur{خالی}}
\textbf{\ur{خالی}} — \emph{khālī} — \textquotedblleft{}empty\textquotedblright{}.

\begin{cousinweb}
\textbf{\ur{خالی}} (\emph{khālī}) is a loan from Arabic \emph{khālī}, \textquotedblleft{}empty\textquotedblright{}, which reached Urdu through Persian. Its Arabic root is \emph{kh-l-w} (Lane).
\end{cousinweb}

\begin{grammarlens}[title={What you've built}]
One more word for this chapter.
\end{grammarlens}

\subsection*{Your turn}
\begin{itemize}
\raggedright
  \item \emph{Say it:} \emph{khālī}
  \item \emph{Say it:} \emph{khālī}, once more
  \item \emph{Say it:} say \emph{khālī}
  \item \emph{Recall:} say the Urdu for low, then the Urdu for full, then say \emph{khālī} again
\end{itemize}

\begin{tcolorbox}[breakable,colback=teal!4,colframe=teal!35!black,title={Before you move on}]
What does \ur{خالی} mean? (\textquotedblleft{}Empty\textquotedblright{}.) Say it once more.
\end{tcolorbox}
\section[\texorpdfstring{\ur{میٹھا} (mīṭhā)}{میٹھا (mīṭhā)}]{\ur{میٹھا} (mīṭhā) — sweet}

\label{lesson:UR-C72-mitha}



\begin{quote}
Before the new one: say the Urdu for full, then the Urdu for empty.
\end{quote}

\subsection*{You'll want to know: \ur{میٹھا}}
\textbf{\ur{میٹھا}} — \emph{mīṭhā} — \textquotedblleft{}sweet\textquotedblright{}.

\begin{cousinweb}
\textbf{\ur{میٹھا}} (\emph{mīṭhā}) is a tadbhava, a word inherited from Sanskrit and reshaped by sound change. It comes from \emph{miṣṭa}, \textquotedblleft{}dainty, sweet\textquotedblright{} (CDIAL 10299). CDIAL lists it under \emph{mṛṣṭa}, \textquotedblleft{}pure; prepared, savoury\textquotedblright{}, of which \emph{miṣṭa} is an early variant.
\end{cousinweb}

\begin{grammarlens}[title={What you've built}]
One more word for this chapter.
\end{grammarlens}

\subsection*{Your turn}
\begin{itemize}
\raggedright
  \item \emph{Say it:} \emph{mīṭhā}
  \item \emph{Say it:} \emph{mīṭhā}, once more
  \item \emph{Say it:} say \emph{mīṭhā}
  \item \emph{Recall:} say the Urdu for full, then the Urdu for empty, then say \emph{mīṭhā} again
\end{itemize}

\begin{tcolorbox}[breakable,colback=teal!4,colframe=teal!35!black,title={Before you move on}]
What does \ur{میٹھا} mean? (\textquotedblleft{}Sweet\textquotedblright{}.) Say it once more.
\end{tcolorbox}
